\documentclass[12pt, a4paper, oneside]{article}
\usepackage[T1]{fontenc}
\usepackage{amsmath, amsthm, amssymb, bm, bookmark, booktabs, color, enumitem, float, graphicx, hyperref, mathrsfs, titling}
\usepackage[UTF8, scheme = plain, fontset = none]{ctex}
\usepackage{listings, subfigure}
\usepackage{grffile}
\title{\textbf{Homework 3}}
\setlength{\droptitle}{-10em}
\author{PENG Qiheng \\ Student ID\: 225040065}
\date{\today}
\linespread{1.5}
\newcounter{problemname}
\newenvironment{problem}{\stepcounter{problemname}\par\noindent{Problem \arabic{problemname}. }}{\par}
\newenvironment{solution}{\par\noindent{Solution. }}{\par}

\begin{document}

\maketitle

\begin{problem}
\begin{enumerate}[label = (\alph*)]
    \item Prove (using words) whether the proposition
    \[
        ((a \lor b) \land (\neg a \lor c)) \to \neg(b \lor c)
    \]
    is valid.

    \begin{solution}
    This proposition is \textbf{not} valid (not a tautology).  
    Give a counterexample:
    \[
        a = \text{F},\quad b = \text{T},\quad c = \text{F}.
    \]
    Then
    \[
        (a \lor b) = \text{T},\quad (\neg a \lor c)=\text{T},
    \]
    so the antecedent is true. But
    \[
        \neg(b \lor c) = \neg(\text{T} \lor \text{F})=\neg\text{T}=\text{F}.
    \]
    Therefore ``true $\to$ false'' is false. So the proposition is not valid.
    \end{solution}

    \item Use propositional resolution to show that the clause set
    \[
    \{p,q,\neg r,s\},\ \{\neg p,r,s\},\ \{\neg q,\neg r\},\ \{p,\neg s\},\ \{\neg p,\neg r\},\ \{r\}
    \]
    is unsatisfiable.

    \begin{solution}
    Let the clauses be
    \[
    \begin{aligned}
    C_1&=(p \lor q \lor \neg r \lor s), \\
    C_2&=(\neg p \lor r \lor s), \\
    C_3&=(\neg q \lor \neg r), \\
    C_4&=(p \lor \neg s), \\
    C_5&=(\neg p \lor \neg r), \\
    C_6&=(r).
    \end{aligned}
    \]
    Resolve step by step:
    \[
    \begin{aligned}
    C_3,\ C_6 &\vdash (\neg q) \qquad\quad\ \ \ \ \ \ \ \ \ \ \ \ \ \ \ \ \ \, (R_1)\\
    C_5,\ C_6 &\vdash (\neg p) \qquad\quad\ \ \ \ \ \ \ \ \ \ \ \ \ \ \ \ \ \, (R_2)\\
    C_4,\ R_2 &\vdash (\neg s) \qquad\quad\ \ \ \ \ \ \ \ \ \ \ \ \ \ \ \ \ \, (R_3)\\
    C_1,\ C_6 &\vdash (p \lor q \lor s) \qquad\ \ \ \ \ \ \ \ \ \ \ \ \ \ \ \, (R_4)\\
    R_4,\ R_1 &\vdash (p \lor s) \qquad\quad\ \ \ \ \ \ \ \ \ \ \ \ \ \ \ \ \ \, (R_5)\\
    R_5,\ R_2 &\vdash (s) \qquad\qquad\quad\ \ \ \ \ \ \ \ \ \ \ \ \ \ \ \ \ \, (R_6)\\
    R_6,\ R_3 &\vdash \Box
    \end{aligned}
    \]
    We derived the empty clause $\Box$, so the clause set is unsatisfiable.
    \end{solution}
\end{enumerate}
\end{problem}

\newpage
\begin{problem}
\begin{enumerate}[label = (\alph*)]
    \item How many of the disjunctions
    \[
    p \lor \neg q,\quad \neg p \lor q,\quad q \lor r,\quad q \lor \neg r,\quad \neg q \lor \neg r
    \]
    can be made simultaneously true?

    \begin{solution}
    All \textbf{five} can be true simultaneously.  
    One satisfying assignment is
    \[
        p=\text{T},\quad q=\text{T},\quad r=\text{F}.
    \]
    Check:
    \[
    \begin{aligned}
    p\lor\neg q &= \text{T}\lor\text{F}=\text{T},\\
    \neg p\lor q &= \text{F}\lor\text{T}=\text{T},\\
    q\lor r &= \text{T}\lor\text{F}=\text{T},\\
    q\lor\neg r &= \text{T}\lor\text{T}=\text{T},\\
    \neg q\lor\neg r &= \text{F}\lor\text{T}=\text{T}.
    \end{aligned}
    \]
    So the maximum number is $5$.
    \end{solution}

    \item In standard mathematics, can
    \[
        x=y,\quad y=z,\quad z=-x
    \]
    hold simultaneously?  
    In propositional logic, consider
    \[
        A\Rightarrow B,\quad B\Rightarrow C,\quad C\Rightarrow \neg A.
    \]
    Decide:
    \begin{enumerate}[label=\roman*.]
        \item They cannot be true simultaneously in all cases?
        \item They can be true simultaneously in some cases?
        \item They can be true simultaneously in all cases?
    \end{enumerate}

    \begin{solution}
    \textbf{Mathematics part.}  
    From $x=y$ and $y=z$, we get $x=z$. Together with $z=-x$, we get $x=-x$, hence $2x=0$, so $x=0$. Then $y=z=0$.  
    Therefore they \textbf{can} hold simultaneously, but only for
    \[
        x=y=z=0.
    \]

    \textbf{Propositional logic part.}  
    Let
    \[
        F=(A\Rightarrow B)\land(B\Rightarrow C)\land(C\Rightarrow\neg A).
    \]
    Truth table:
    \[
    \begin{array}{ccccccc}
    A & B & C & A\Rightarrow B & B\Rightarrow C & C\Rightarrow\neg A & F\\
    \hline
    T & T & T & T & T & F & F\\
    T & T & F & T & F & T & F\\
    T & F & T & F & T & F & F\\
    T & F & F & F & T & T & F\\
    F & T & T & T & T & T & T\\
    F & T & F & T & F & T & F\\
    F & F & T & T & T & T & T\\
    F & F & F & T & T & T & T
    \end{array}
    \]
    Conclusions:
    \begin{enumerate}[label=\roman*.]
        \item Yes. They are \textbf{not} true in all cases (for example $A=B=C=T$ makes $F$ false).
        \item Yes. They are true in some cases (for example $A=B=C=F$ makes $F$ true).
        \item No. They are not true in all cases.
    \end{enumerate}
    \end{solution}
\end{enumerate}
\end{problem}

\end{document}
